\documentclass[11pt,a4paper]{article}

\usepackage[T1]{fontenc}
\usepackage[utf8]{inputenc}
\usepackage[margin=1.05in]{geometry}
\usepackage{amsmath,amssymb,amsthm,mathtools}
\usepackage{booktabs}
\usepackage{graphicx}
\usepackage[protrusion=true,expansion=false]{microtype}
\usepackage{siunitx}
\usepackage{caption}
\usepackage{enumitem}
\usepackage[table]{xcolor}
\usepackage[hidelinks,breaklinks]{hyperref}
\usepackage{cleveref}

\sisetup{detect-weight=true, group-separator={,}}

\theoremstyle{definition}
\newtheorem{definition}{Definition}[section]
\theoremstyle{plain}
\newtheorem{proposition}[definition]{Proposition}
\theoremstyle{remark}
\newtheorem{remark}[definition]{Remark}
\newtheorem{assumption}[definition]{Assumption}

\newcommand{\R}{\mathbb{R}}
\newcommand{\E}{\mathbb{E}}
\newcommand{\rank}{\operatorname{rank}}
\newcommand{\VOLL}{\mathrm{VOLL}}
\newcommand{\MW}{\,\mathrm{MW}}
\newcommand{\MWh}{\,\mathrm{MWh}}
\newcommand{\okina}{\textquoteleft}
\newcommand{\Oahu}{O\okina ahu}

\definecolor{filed}{HTML}{1E4E8C}
\definecolor{assum}{HTML}{B86F12}
\newcommand{\Filed}{\textcolor{filed}{\textsc{filed}}}
\newcommand{\Assum}{\textcolor{assum}{\textsc{assumed}}}

\captionsetup{font=small,labelfont=bf}

\title{\bfseries Machines, Not Dials:\\[0.25em]
\large Chronological Unit Commitment and a Reconstructed\\
Spatial Demand Field for the \Oahu{} Power System}

\author{Pascual Z\'arate\thanks{The University of Chicago, Class of 2029.
Prepared as Set 9 of an ongoing research project on island power-system
optimisation. All code, data and numerical results are reproducible from the
repository described in \Cref{app:repro}.}}

\date{September 2026}

\begin{document}
\maketitle

\begin{abstract}
\noindent
A companion study (Set 8) found that the distributional assumption placed on
nodal demand relocates, rather than merely inflates, estimated adequacy risk in
a DC optimal power flow model of the \Oahu{} grid. That study represented
generators by a single box constraint, $0 \le P_g \le P_g^{\max}$, and sampled
hours independently with replacement. Both choices are load-bearing: the first
removes every operating limit a thermal machine actually has, and the second
destroys the temporal adjacency that such limits act upon. This paper rebuilds
the experiment as a chronological mixed-integer unit commitment problem over
$8{,}760$ consecutive hours, adding minimum stable generation, ramp limits,
minimum up and down times, and start-up costs one constraint at a time.

Two results follow. First, the binding constraint on this system is
\emph{ramp}, not capacity: at a reserve margin of $43\%$ the fleet sheds energy
in $0.61\%$ of hours once hour-to-hour movement is restricted, and in no hours
before that. The shortfall is one of speed, and the earlier model was
structurally blind to it. Second, and more consequentially for method, we show
that the spatial demand field used by Set 8 --- and by most studies that have
only a system-level load series --- is \emph{rank one}: a fixed share vector
scaled by a scalar time series. Its direction cannot move, yet direction is the
quantity the research question is about. We reconstruct a rank-two field from
three public datasets, preserving the filed system total to
$3\times10^{-4}\MW$ while allowing the share of load at the urban node to swing
from $52.7\%$ to $84.1\%$ across the year, a diurnal rotation that emerges
without being imposed.

We report two errors found by self-audit during this work, with their
quantified effect, and a $600$-run sensitivity protocol that converts the
principal remaining judgement call into a measured statement. We conclude that
every \emph{directional} result here is robust, and that no network
\emph{magnitude} should be read as a claim about the physical \Oahu{} system,
because the network representation is entirely assumed.
\end{abstract}

\section{Introduction}
\label{sec:intro}

Island power systems are the cleanest available laboratory for adequacy
questions. \Oahu{} has no interconnection: if the island cannot generate the
energy, no neighbour can supply it. Every megawatt-hour served is produced
within a $1{,}508\MW$ fleet of $24$ thermal units at six sites.

Set 8 of this project asked what happens to computed adequacy risk when the
analyst's distributional assumption about demand is wrong. The answer was that
misspecification does not merely scale the risk estimate; it moves the risk to
different nodes. That finding rested on a dispatch model in which a generator
was a single interval constraint. A $128\MW$ steam unit could idle at $3\MW$,
reach full output within the hour, and cycle on and off at no cost.

This is not a small idealisation, and it is not obviously conservative. It is
not clear \emph{a priori} whether removing operating limits makes a system look
safer (any unit can always be called upon) or more dangerous (no unit is ever
forced to stay on). The present paper resolves that question empirically and
finds something more interesting than either: the limits create an entirely new
failure mode, invisible to the earlier model, that has nothing to do with
installed capacity.

\paragraph{Contributions.}
\begin{enumerate}[leftmargin=1.4em,itemsep=2pt]
\item A chronological MILP unit commitment formulation over the full year
  (\Cref{sec:model}), with an explicit treatment of the start-up ramp that is
  \emph{necessary} rather than refining --- omitting it renders every steam unit
  on this fleet permanently unstartable (\Cref{sec:startup}).
\item A reproducible reconstruction of a spatially and temporally resolved
  demand field from public data, raising the demand matrix from rank one to
  rank two while preserving the filed system total exactly
  (\Cref{sec:demand}).
\item A constraint ladder that attributes every movement in the results to a
  single added restriction (\Cref{sec:ladder,sec:results}).
\item A self-audit reporting two errors in our own construction with their
  quantified effect, and a sensitivity protocol over $600$ perturbed model
  variants (\Cref{sec:audit,sec:sensitivity}).
\end{enumerate}

\paragraph{A note on epistemic labelling.}
Throughout, quantities are marked \Filed{} when they originate in a government
form or a regulatory filing, and \Assum{} when they are our own judgement. This
is not decoration. \Cref{sec:limits} shows that the two findings of this paper
rest on layers of very different evidential quality, and the distinction
determines which of them may be quoted as a statement about \Oahu{}.

\section{The \Oahu{} test system}
\label{sec:system}

\subsection{Generation}

The fleet is taken from EIA Form 860 (2024, summer capacity) restricted to
utility and contracted firm \Oahu{} plants: Kahe, Waiau, H-POWER, Kalaeloa,
Campbell Industrial Park, and Schofield Generating Station. Airport emergency
generators and refinery self-generation are excluded, as are solar, wind and
battery resources, whose effect is already embedded in the net load series.
Kalaeloa's two combustion turbines and steam tail are aggregated as a single
$220\MW$ combined-cycle resource; dispatching the steam tail independently
would be physically misleading.

\begin{table}[htbp]
\centering
\caption{Composition of the modelled fleet. Capacities are \Filed{};
operating limits are \Assum{} by technology class.}
\label{tab:fleet}
\small
\begin{tabular}{lrrrr}
\toprule
Technology & Units & Capacity (\si{\mega\watt}) & $P^{\min}/P^{\max}$ & Ramp (\%/h) \\
\midrule
Oil steam            & 12 & 936.3 & 0.40 & 30 \\
Combined cycle       &  1 & 220.0 & 0.40 & 50 \\
Combustion turbine   &  3 & 215.4 & 0.00 & 100 \\
Municipal waste      &  2 &  86.0 & 0.80 & 20 \\
Internal combustion  &  6 &  50.4 & 0.00 & 100 \\
\midrule
\textbf{Total}       & \textbf{24} & \textbf{1{,}508.1} & & \\
\bottomrule
\end{tabular}
\end{table}

Half the fleet by unit count, and $62.1\%$ of it by capacity, is oil steam ---
the class combining a high minimum stable generation with the slowest ramp. The
genuinely flexible plant (combustion turbines and reciprocating engines, which
may start from cold and move at full rate) totals only $265.8\MW$, or $17.6\%$
of capacity. This structural fact drives the principal result of
\Cref{sec:results}.

\subsection{Demand}

The load series is the 2021 base scenario of Hawaiian Electric's Integrated
Grid Planning \Oahu{} Inputs Workbook 3, filed 2022-03-31 in PUC docket
2018-0165. It is assembled from ten hourly layers as
\[
  \text{gross} = \text{underlying} + \text{managed EV} - \text{base forecast}
                 + \text{future eBus},
\]
with net load subtracting distributed PV, four storage sectors, efficiency and
time-of-use shifting. The result is $8{,}760$ observations with mean
$743.9\MW$, minimum $472.6\MW$ and peak $1{,}054.3\MW$.

\begin{remark}[Filed is not metered]
\label{rem:filed}
This series is a \emph{filed planning profile}, not a meter reading. It was
submitted under oath to the regulator, so it is a real document; but it
records what the utility modelled, not what the system recorded. We use
``filed'' rather than ``measured'' throughout, and recommend the same
discipline to anyone quoting these results.
\end{remark}

The reserve margin implied by \Cref{tab:fleet} against the annual peak is
\[
  \frac{1508.1 - 1054.3}{1054.3} = 43.0\%.
\]
This number should be held in mind throughout \Cref{sec:results}: the system is
not short of installed capacity.

\subsection{Network}
\label{sec:network}

The network is a stylised nine-bus representation: the seven \Oahu{} judicial
districts as published by the Census as county subdivisions, with the Ewa
district split into \textsc{Ewa-West}, \textsc{Ewa-Central} and \textsc{Kahe}.
The split is necessary: point-in-polygon assignment on EIA-860 coordinates
places all six modelled plants inside Ewa, which spans roughly twenty miles, so
leaving it whole would collapse all island generation onto one node and erase
precisely the injection geography the study exists to examine. Ten branches are
wired as the two-corridor ring (one north, one south) that Hawaiian Electric
describes in public material, leaving two independent loops --- a requirement,
since a radial network makes DC-OPF and the transportation model identical.

\begin{assumption}[The network is ours]
\label{as:network}
Bus \emph{locations} are \Filed{}: population-weighted 2020 census-tract
centroids, and EIA plant coordinates. Every \emph{electrical} parameter is
\Assum{}: the topology, line lengths (great-circle separation $\times 1.25$),
reactance ($0.80\,\Omega$/mi at \SI{138}{\kilo\volt} on a \SI{100}{\mega\volt\ampere}
base), thermal limits ($200\MW$ per circuit) and circuit counts. Not one comes
from a filing.
\end{assumption}

\Cref{as:network} is the reason \Cref{sec:congestion} declines to interpret the
magnitude of any congestion statistic.

\section{Chronological unit commitment}
\label{sec:model}

\subsection{Formulation}

Let $\mathcal{G}$ index units, $\mathcal{N}$ buses, $\mathcal{L}$ branches and
$\mathcal{T}$ the hours of a block. Binary $u_{g,t}$ denotes commitment, with
start-up $v_{g,t}$ and shut-down $w_{g,t}$. Continuous $P_{g,t}$ is output,
$\theta_{n,t}$ the voltage angle and $s_{n,t} \ge 0$ load shed. The objective is
\begin{equation}
\min \sum_{t\in\mathcal{T}}\Bigl[
  \sum_{g\in\mathcal{G}} \bigl(c_g P_{g,t} + C^{\mathrm{su}}_g v_{g,t}\bigr)
  + \VOLL \sum_{n\in\mathcal{N}} s_{n,t}
  + \epsilon \sum_{g} g\,P_{g,t}
\Bigr],
\label{eq:obj}
\end{equation}
with $\VOLL = \SI{10000}{\text{USD}\per\mega\watt\hour}$ and a tie-break term
$\epsilon = 10^{-6}$ that makes dispatch unique between cost-identical units.
The tie-break is a correctness device, not a cosmetic one: without it, equal-cost
units receive an arbitrary split and results are not reproducible across solver
versions.

Subject to nodal balance, DC flow $f_{\ell,t} = b_\ell(\theta_{n,t}-\theta_{m,t})$
with $|f_{\ell,t}| \le \bar f_\ell$, an angular reference at the urban bus, and
the operating constraints
\begin{align}
P^{\min}_g u_{g,t} &\le P_{g,t} \le P^{\max}_g u_{g,t}, \label{eq:pmin}\\
u_{g,t} - u_{g,t-1} &= v_{g,t} - w_{g,t}, \label{eq:logic}\\
\textstyle\sum_{\tau=t-T^{\mathrm{up}}_g+1}^{t} v_{g,\tau} &\le u_{g,t},
\qquad
\textstyle\sum_{\tau=t-T^{\mathrm{dn}}_g+1}^{t} w_{g,\tau} \le 1-u_{g,t}.
\label{eq:minupdn}
\end{align}

\subsection{The start-up ramp is necessary, not refining}
\label{sec:startup}

The naive ramp constraint $|P_{g,t}-P_{g,t-1}| \le R_g$ is wrong here, and
catastrophically so. On this fleet every steam unit has
$P^{\min}_g = 0.40 P^{\max}_g$ but $R_g = 0.30 P^{\max}_g$. Starting up requires
moving from $0$ to $P^{\min}_g$ within one hour by \eqref{eq:pmin}, a step of
$0.40P^{\max}_g > R_g$. The unit is therefore \emph{permanently unstartable}.

\begin{proposition}
If $P^{\min}_g > R_g$ and ramping is enforced without a start-up allowance, then
any unit off at $t=0$ satisfies $u_{g,t}=0$ for all $t$.
\end{proposition}

Our first implementation did exactly this and reported load shedding in
$95.8\%$ of hours --- an apparent collapse of the \Oahu{} system that was
entirely an artefact of formulation. The correct form separates the start-up
movement, with $S_g = \max(P^{\min}_g, R_g)$:
\begin{align}
P_{g,t} - P_{g,t-1} &\le R_g u_{g,t-1} + S_g v_{g,t}, \label{eq:rampup}\\
P_{g,t-1} - P_{g,t} &\le R_g u_{g,t}   + S_g w_{g,t}. \label{eq:rampdn}
\end{align}

\begin{remark}[The diagnostic was economic, not numerical]
\label{rem:tell}
The solver reported optimality at every block. No exception was raised. The
error surfaced because the \emph{dispatch pattern} was economically absurd:
cheap steam capacity idle while expensive diesel engines ran continuously, with
$43\%$ of capacity unused during a purported shortage. We record this because
it generalises --- in mixed-integer dispatch models, implausible economics is a
more sensitive detector of specification error than any solver diagnostic.
\end{remark}

\subsection{Rolling horizon}

Ramp constraints couple $t$ to $t-1$; minimum up/down times couple a commitment
decision to several subsequent hours. Set 8 sampled hours independently with
replacement, so hour $5{,}000$ could follow hour $12$ and no ``previous hour''
existed for \eqref{eq:rampup} to reference. Set 9 therefore solves the year as
$365$ consecutive $24$-hour blocks, carrying $(u,P,T^{\mathrm{up}},T^{\mathrm{dn}})$
across each boundary. This is a different experimental design, not a refinement,
which is why Set 9 reuses none of Set 8's numerical results. Blocks are capped
at a $1\%$ optimality gap and $20$ seconds; capped blocks are counted and
reported.

\section{Reconstructing the spatial demand field}
\label{sec:demand}

\subsection{The rank-one problem}

Let $D \in \R^{9 \times 8760}$ be the bus-by-hour demand matrix. Set 8, like
essentially every study with only a system-level load series, constructed
\begin{equation}
  D = \alpha\, \ell^{\top}, \qquad \alpha \in \Delta^{8}, \;\; \ell \in \R^{8760},
\label{eq:rank1}
\end{equation}
with $\alpha$ the 2020 resident-population share vector. Then
$\rank(D) = 1$: the $78{,}840$ entries carry only $8{,}760$ degrees of freedom,
the nine bus loads can only scale together, and the demand \emph{direction}
$D_{:,t}/\|D_{:,t}\|_1$ is constant in $t$.

This is a serious problem for the present research question. Set 8's own
finding is that direction governs adequacy. Equation \eqref{eq:rank1} holds
direction fixed by construction, so the quantity of interest was the one
quantity the data could not express.

\subsection{Static sector reallocation}

The first repair replaces population with a sector-weighted key:
\begin{equation}
  \alpha_i = f_{\mathrm{res}}\frac{\mathrm{pop}_i}{\sum_j \mathrm{pop}_j}
           + f_{\mathrm{com}}\frac{\mathrm{com}_i}{\sum_j \mathrm{com}_j}
           + f_{\mathrm{ind}}\frac{\mathrm{ind}_i}{\sum_j \mathrm{ind}_j},
\label{eq:v2}
\end{equation}
with sector energy shares $f$ from EIA Form 861 restricted to Hawaiian Electric
Co Inc, the \Oahu{} operating entity ($f_{\mathrm{res}},f_{\mathrm{com}},
f_{\mathrm{ind}} = 0.2675, 0.3112, 0.4213$), residential sited by Census
population and the two employment sectors by LEHD LODES workplace jobs.
Normalising each key separately is what makes the terms commensurable; the
shares then sum to unity automatically, since each key sums to one and
$\sum_s f_s = 1$.

\Oahu{} is a commuter island: the Honolulu district holds $39.9\%$ of residents
but $71.0\%$ of jobs. Allocating by residence therefore exports a large block of
commercial load out of the urban core in every hour. Under \eqref{eq:v2} the
Honolulu share rises from $39.87\%$ to $56.89\%$. The method is not simply
urban-biased: \textsc{Ewa-West} \emph{gains} share ($2.96\% \to 4.73\%$) despite
few residents, because Campbell Industrial Park places $7.76\%$ of the island's
industrial employment there against $2.96\%$ of its population.

But \eqref{eq:v2} is still of the form \eqref{eq:rank1}. The rank is still one.

\subsection{Sector-resolved hourly allocation}
\label{sec:v3}

The rank barrier is structural: any single share vector, however well
estimated, yields a rank-one field. Rank rises only if sectors with
\emph{different spatial footprints} also have \emph{different temporal
profiles}. Writing $a_{i,s}$ for the share of sector $s$ at bus $i$ and
$\gamma_s(t)$ for its system-level hourly load,
\begin{equation}
  D_{i,t} = \sum_{s} a_{i,s}\,\gamma_s(t) \;-\; \pi_i\, \mathrm{dg}(t),
\label{eq:v3}
\end{equation}
where $\mathrm{dg}(t)$ is the behind-the-meter layer and $\pi$ its (assumed)
spatial key. Sector profiles $\gamma_s$ are built from OpenEI TMY3 reference
buildings at Honolulu weather, folded into $12 \times 2 \times 24 = 576$
month/weekend/hour cells so that repeatable rhythm is retained and single-year
weather noise is averaged out. The commercial profile is an energy-weighted
composite of fifteen reference building types.

Two implementation details matter. Behind-the-meter PV can exceed a suburban
bus's gross draw at midday, which \eqref{eq:v3} would render negative; since the
network model admits no negative injection at a load bus, such values are
clipped at zero and the surplus redistributed proportionally. This is a genuine
loss of realism --- \Oahu{} does have circuits exporting at midday --- and it is
visible in the results as a minimum share of exactly $0.00\%$ at
\textsc{Waianae}. Finally, all bus loads are rescaled so that
$\sum_i D_{i,t}$ reproduces the filed net load exactly.

\subsection{Verification}

\begin{table}[htbp]
\centering
\caption{Verification of the reconstructed field. Only the distribution of
demand is altered; its total is preserved.}
\label{tab:verify}
\small
\begin{tabular}{lrrr}
\toprule
 & v1 (population) & v2 (sector) & v3 (hourly) \\
\midrule
Matrix rank                       & 1 & 1 & \textbf{2} \\
Singular values                   & --- & --- & 62.7, 5.79, 0.14 \\
Honolulu share, minimum           & 39.87\% & 56.89\% & 52.71\% \\
Honolulu share, maximum           & 39.87\% & 56.89\% & \textbf{84.10\%} \\
Annual swing (points)             & 0.00 & 0.00 & \textbf{31.38} \\
Max error vs filed total (\si{\mega\watt}) & --- & --- & $3\times10^{-4}$ \\
\bottomrule
\end{tabular}
\end{table}

The rank is two rather than three because industrial load is modelled flat and
therefore contributes no independent temporal variation; the second dimension
is entirely residential-versus-commercial. We regard reporting rank two as more
honest than engineering a third from data we do not have.

The diurnal pattern is not imposed. Averaged over the year, the Honolulu share
runs $55.8\%$ at midnight, $75.4\%$ at noon and $53.7\%$ at 22:00. Two
mechanisms push the same way at midday: commercial load peaks in the urban core
while distributed PV suppresses net load in the residential periphery.

\section{Experimental design}
\label{sec:ladder}

Constraints are added one at a time, so that any movement in a reported
statistic is attributable to a single restriction:
\begin{enumerate}[label=(\roman*),leftmargin=2.2em,itemsep=1pt]
\item \texttt{lp}: $0 \le P \le P^{\max}$ --- Set 8's constraint set, as control;
\item \texttt{pmin}: binary commitment and \eqref{eq:pmin};
\item \texttt{ramp}: \eqref{eq:rampup}--\eqref{eq:rampdn};
\item \texttt{full}: \eqref{eq:minupdn} and start-up costs.
\end{enumerate}

\section{Results}
\label{sec:results}

\subsection{The constraint ladder}

\begin{table}[htbp]
\centering
\caption{The constraint ladder under the reconstructed (v3) demand field, full
year. ENS is energy not served.}
\label{tab:ladder}
\small
\begin{tabular}{lrrrr}
\toprule
Rung & Hours at a line limit & Hours shedding & ENS (\si{\mega\watt\hour}) & Cost (\$M) \\
\midrule
\texttt{lp}   & 94.77\% & 0.00\% & 0.0  & 688.5 \\
\texttt{pmin} & 58.93\% & 0.00\% & 0.0  & 693.3 \\
\texttt{ramp} & 90.37\% & \textbf{0.61\%} & \textbf{33.8} & 689.0 \\
\texttt{full} & 62.66\% & 0.21\% & 11.7 & 689.8 \\
\bottomrule
\end{tabular}
\end{table}

\subsection{Ramp, not capacity, is the binding constraint}

The system serves all load at the first two rungs. Energy is shed only once
\eqref{eq:rampup} is imposed --- in $0.61\%$ of hours, about $53$ hours per
year, totalling $33.8\MWh$. Nothing about installed capacity changes between
\texttt{pmin} and \texttt{ramp}; only the permitted rate of change does. At a
$43\%$ reserve margin, the machines exist and cannot arrive in time.

This is the paper's substantive new result, and it is a category of risk the
earlier model could not represent, because a model without temporal adjacency
has nothing for a rate limit to bind against. It is also small: $33.8\MWh$
against annual consumption near $6.5\,\mathrm{TWh}$. It is not zero, and
storage is precisely the technology that removes it, which makes storage the
natural next study rather than a generic recommendation.

Shedding \emph{falls} from \texttt{ramp} to \texttt{full}. This is not
paradoxical: minimum up-times force units to remain committed through quiet
hours, so more capacity is already synchronised and able to move when demand
rises. Commitment inertia is a cost in fuel and a benefit in flexibility.

\subsection{Congestion: direction without magnitude}
\label{sec:congestion}

The fraction of hours with at least one line at its limit moves substantially
across the ladder ($94.77\% \to 58.93\% \to 90.37\% \to 62.66\%$). The pattern
is interpretable --- committing fewer, larger units concentrates injection and
loads the corridors; minimum up-times then spread it again --- but by
\Cref{as:network} every line rating and reactance producing these numbers is
our own. We therefore report the direction and explicitly decline to interpret
the magnitude. These are statements about a nine-bus sketch, not about the
\Oahu{} transmission system.

\section{Self-audit}
\label{sec:audit}

Two errors in our own construction were found after the first complete set of
results, and both changed the answer. We report them with their effect.

\subsection{Reference-calendar misalignment}

The OpenEI reference-building simulations run on a calendar year beginning on a
Sunday; 2021 begins on a Friday. Both are non-leap, so \emph{dates} correspond
exactly and our implementation assumed \emph{weekdays} did too. They are two
days apart. The consequence was that simulated Saturday and Sunday were filed
as Thursday and Friday: a large office read as busier at the weekend than
midweek. The diagnostic signature is unmistakable in the raw profile --- the
2\,p.m.\ load for a large office in June reads $1{,}568$, $1{,}568$, $953$,
$532$, $1{,}566$, $1{,}567$, $1{,}567$ \si{\kilo\watt} on days labelled Tuesday
through Monday. The weekly pattern is correct; it sits on the wrong days.

The offset was confirmed independently against three building types with a
genuine work-week signal (large office, primary school, hospital), all implying
$+2$. Correcting it moved the weekend-to-weekday commercial ratio from $1.14$ to
$0.63$ at the profile level. The hour-of-day curve moved by at most $0.036$ of
its peak and the peak hour remained 13:00 --- so the headline diurnal result
was unaffected, while the weekly rhythm had been inverted.

\subsection{Headcount as a proxy for energy}

The commercial composite was initially weighted by employment alone, which
assumes an office worker and a hotel worker represent equal load. EIA CBECS
2018 (table C22 over table B1) gives electricity per worker by principal
building activity: $58{,}888$ \si{\kilo\watt\hour} for food sales, $36{,}179$
for lodging, $20{,}266$ for mercantile, and $6{,}912$ for office --- the lowest
of any activity measured. Re-weighting by estimated energy moves offices from
$39.5\%$ of the commercial blend to $18.4\%$, with lodging/food and retail each
gaining roughly twelve points.

The substantive finding is that \Oahu{}'s commercial load is
\emph{tourism-shaped rather than office-shaped}, and headcount weighting
concealed it. The observable consequence is again the working week: hotels,
restaurants and shops trade seven days where offices do not, so the
weekend-to-weekday ratio rises from $0.757$ to $0.809$ in the delivered field.
CBECS intensities are national; applied to a tourism economy they likely
understate lodging, making this correction conservative.

\section{Sensitivity}
\label{sec:sensitivity}

The allocation of each employment category across building types
(\Cref{sec:v3}) has no source and remains our judgement. Rather than declare it,
we tested it: the commercial profile was rebuilt under $300$ random
within-category mixes and $300$ random between-category re-weightings spanning
a ninefold range, measuring the maximum deviation of the normalised hour-of-day
curve.

\begin{table}[htbp]
\centering
\caption{Sensitivity of the hour-of-day commercial profile. A deviation of
$1.000$ would indicate complete change of shape.}
\label{tab:sens}
\small
\begin{tabular}{lrr}
\toprule
Perturbation & Max deviation & Peak hour \\
\midrule
Equal weights within category      & 0.007 & 13:00 \\
Largest building type only         & 0.055 & 13:00 \\
Health care $\times 3$             & 0.070 & 13:00 \\
Lodging and food $\times 3$        & 0.090 & 13:00 \\
\textbf{Worst of 600 random variants} & \textbf{0.104} & \textbf{13:00} \\
\bottomrule
\end{tabular}
\end{table}

The peak hour never moved. The mechanism is transparent: every commercial
building type shares a daytime-peaking profile, and averaging mutually
consistent curves cannot produce much disagreement. This contrasts instructively
with the calendar error of \Cref{sec:audit}, which was \emph{structural} --- it
relocated load onto the wrong days --- and therefore mattered far more than a
large error in weights.

The result's own boundary is equally clear. It holds because nothing in the
average runs at night. Uniformed military personnel are excluded from LODES, so
Pearl Harbor--Hickam, Schofield Barracks and K\=ane\okina ohe Marine Corps Base
--- large, continuously operating loads --- contribute nothing to these
profiles. That, not the building mix, is the live exposure.

\section{Limitations}
\label{sec:limits}

\begin{enumerate}[leftmargin=1.4em,itemsep=2pt]
\item \textbf{Operating limits are assumed.} Neither EIA-860 nor EIA-923 carries
  unit-level minimum stable output or ramp rate, and the IGP workbooks expose no
  traceable unit table for them. \Cref{tab:fleet} therefore rests on technology
  class conventions. The ramp result of \Cref{sec:results} depends entirely on
  these values; obtaining verified operating data is the highest-value
  improvement available to this model.
\item \textbf{The network is assumed} (\Cref{as:network}).
\item \textbf{Military load is absent}, as above.
\item \textbf{Vintage mismatch.} The fleet is 2024 vintage; the load series is
  2021. AES Hawaii ($180\MW$ of coal) retired in September 2022 and is correctly
  absent from the fleet, but served the 2021 load. The model therefore asks a
  post-AES fleet to serve pre-AES demand --- a defensible scenario, but not a
  reconstruction of either year, and it should be named as such.
\item \textbf{No forced outages, no $N-1$, no storage.} Every unit is available
  in every hour. For an adequacy study this remains the largest single gap.
\item \textbf{Reference-building profiles} are typical-year simulations, not
  \Oahu{}'s metered sector loads; industrial load is flat; distributed PV is
  allocated by population.
\end{enumerate}

\section{Conclusion}

Modelling \Oahu{}'s generators as machines rather than dials does not overturn
the Set 8 finding that the assumed demand distribution governs computed risk. It
does reveal a failure mode the earlier formulation could not express: at a $43\%$
reserve margin the fleet is unable to follow real demand in roughly $53$ hours
per year, shedding $33.8\MWh$. The constraint that binds is rate, not quantity.

Methodologically, we argue that the rank-one demand field implicit in
system-total load data is a more serious limitation than it is usually treated
as, and that it can be partially repaired from public data. The reconstruction
presented here raises the field to rank two and produces a diurnal rotation of
the demand direction of $31.4$ percentage points, while preserving the filed
system total to $3\times10^{-4}\MW$.

Finally, we note that of the two errors found by self-audit, the damaging one
was structural rather than numerical, while a deliberately mis-specified
building mix --- tested over $600$ variants --- barely moved the answer. We take
this as an argument for auditing the \emph{shape} of an assumption before
refining its parameters.

\appendix

\section{Reproducibility}
\label{app:repro}

All results are reproducible from the project repository. The Set 9 pipeline is:
\begin{enumerate}[leftmargin=1.4em,itemsep=1pt]
\item \texttt{prepare\_demand.py} --- static sector reallocation, \eqref{eq:v2}.
\item \texttt{prepare\_hourly\_demand.py} --- sector-resolved hourly field,
  \eqref{eq:v3}. Both scripts download their own sources.
\item \texttt{fleet.jl} --- case construction; all assumed parameters are
  declared here and nowhere else.
\item \texttt{commitment.jl} --- the MILP, \eqref{eq:obj}--\eqref{eq:rampdn}.
\item \texttt{run\_hourly.jl} --- the ladder under v1, v2 and v3.
\end{enumerate}
Raw archives are SHA-256 pinned in \texttt{SOURCES.md}; the demand
reconstruction, its verification and its sensitivity analysis are documented in
\texttt{DATA\_REBUILD.md}.

\section{Assumed parameters}
\label{app:assumed}

\begin{table}[htbp]
\centering
\small
\caption{Parameters with no documentary source, collected for challenge.}
\begin{tabular}{llr}
\toprule
Parameter & Basis & Value \\
\midrule
$P^{\min}/P^{\max}$, oil steam & technology class & 0.40 \\
Ramp, oil steam & technology class & 30\%/h \\
Minimum up-time, oil steam & convention & 8 h \\
Minimum up-time, municipal waste & continuous refuse feed & 24 h \\
Start-up cost & convention, USD per \si{\mega\watt} capacity & 20--100 \\
Line thermal limit & per \SI{138}{\kilo\volt} circuit & $200\MW$ \\
Line reactance & at \SI{138}{\kilo\volt}, \SI{100}{\mega\volt\ampere} base & $0.80\,\Omega$/mi \\
Routing factor & great-circle multiplier & 1.25 \\
$\VOLL$ & convention & \SI{10000}{\text{USD}\per\mega\watt\hour} \\
Ewa split meridian & by census tract & $-158.075^\circ$ \\
\bottomrule
\end{tabular}
\end{table}

\end{document}
